\section{Notation and exact layer specifications}
\label{app:v4-layers}
\begin{longtable}{p{0.21\textwidth}p{0.72\textwidth}}
\toprule
Symbol & Meaning\\
\midrule
\endhead
$Y$, $\vD$ & Binary target-training association matrix and its unique positive edges.\\
$b_e,b_r$ & Native 512-dimensional phase-1 protein and reaction outputs.\\
$g_e,f_e$ & Separately normalized global and fused protein branches.\\
$p_{ki},a_{ki}$ & Raw and uniform-smoothed residue attention for learned slot $k$.\\
$\omega_{vj}$ & Gate weight for protein view $v$ and hidden coordinate $j$.\\
$s$ & Learned bounded protein fusion scale; initial value 0.1.\\
$\nu$ & Protein fusion multiplier used after training; value 3.\\
$h_E,h_R$ & Independent residual heads fitted with frozen F3 inputs.\\
$\kappa,c$ & Residual coefficient in training (0.2) and inference multiplier (0.5).\\
$u_e^{\mathrm{tr}},u_r^{\mathrm{tr}}$ & Refined endpoints used in the phase-2 training objectives.\\
$u_e,u_r$ & Refined endpoints after the two inference adjustments.\\
$\lambda_f,\xi_c,w_{ic}$ & Family coefficient, category weight, and membership confidence; these have different roles.\\
$a_e,a_r$ & Normalized training-reaction dictionary coordinates.\\
$\Phi_E,\Phi_R$ & Complete $512+N_R$ endpoint representations.\\
$\alpha$ & Dictionary contribution to the final dot product; value 0.4.\\
\bottomrule
\end{longtable}

\begin{longtable}{>{\raggedright\arraybackslash}p{0.30\textwidth}>{\raggedright\arraybackslash}p{0.63\textwidth}}
\caption{Active neural components in the reported configuration.
Dimensions are listed in execution order.}\label{tab:v4-layers}\\
\toprule
Component & Specification\\
\midrule
\endfirsthead
\toprule Component & Specification\\\midrule
\endhead
Frozen residue features & ProtT5; 1024 coordinates per retained residue.\\
Frozen SLEEC scorer & Linear $1024\to256$; ReLU; linear $256\to1$.\\
Residue adapter & LayerNorm(1024); linear $1024\to256$; GELU.\\
Keys and values & Keys: bias-free linear $256\to256$ followed by normalization. Values: affine linear $256\to256$.\\
Learned residue queries & Four trainable 256-dimensional vectors; orthogonal initialization; normalized before scoring.\\
Global/SLEEC view projection & Two independent networks: LayerNorm(1024); linear $1024\to512$; GELU; dropout 0.1; linear $512\to256$; normalization.\\
Learned-slot view projection & Four independent networks: LayerNorm(256); linear $256\to512$; GELU; dropout 0.1; linear $512\to256$; normalization.\\
Protein gate & Linear $1536\to256$; GELU; linear $256\to1536$; reshape to $6\times256$; softmax over views.\\
Global output branch & LayerNorm(1024); linear $1024\to512$; GELU; dropout 0.1; linear $512\to512$; normalization.\\
Fused output branch & LayerNorm(256); linear $256\to512$; GELU; dropout 0.1; linear $512\to512$; normalization.\\
Uni-Mol2 molecule pooling & Separate scalar affine attention for reactants and products, acting on 768-dimensional molecule vectors.\\
ChIRo molecule pooling & Separate scalar affine attention for reactants and products, acting on 256-dimensional molecule vectors.\\
Reaction modality MLPs & Input widths $768,3072,1024,617$ for T5, Uni-Mol2 composition, ChIRo composition, and chemistry. Each follows $d_m\to1024\to512\to512$; ReLU after its first two layers.\\
Reaction token normalization & One LayerNorm(512), shared over modality tokens. No per-token $\ell_2$ normalization in this configuration.\\
Reaction modality attention & Linear $512\to512$; ReLU; linear $512\to1$; masked softmax over available modalities.\\
Reaction output projection & Linear $512\to1024$; ReLU; linear $1024\to512$; normalization.\\
Phase-2 heads & Two independent copies of LayerNorm(512); linear $512\to1024$; GELU; linear $1024\to512$. Last affine layer initialized to zero.\\
Biological neighborhood loss & A nonparametric training helper. Zero additional learned parameters and no inference-time readout.\\
\bottomrule
\end{longtable}

The two phase-2 heads contain 2,102,272 parameter values in total. Biological
supervision changes their gradients but not their parameter names or shapes.
The four residue queries are shared across proteins; they are not four
trainable vectors created separately for each sequence.

\section{Attention implementation and sequence preprocessing}
\label{app:v4-attention}
Padding is removed from softmax denominators and pooling sums, and its
contribution is zero. Before the frozen scorer and residue adapter, padded
values are replaced by zero so that padded NaN sentinels do not enter learned
linear maps. SLEEC logits are detached. Centering its logits alone would not
alter softmax, but centering before clipping gives a protein-relative clipping
rule. The use of logits, rather than sigmoid scores, is material to the
resulting distribution.

The learned attention scale is initialized by
$\rho_\gamma=\operatorname{logit}(4/10)$, and the fusion scalar by
$\rho_s=\operatorname{logit}(0.1/0.5)$. The normalization floor in the
protein multiview module is $10^{-6}$. Attention entropy uses natural
logarithms and a probability floor of $10^{-12}$ inside the logarithm.
Regularization in Eqs.~\eqref{eq:v4-entropy}--\eqref{eq:v4-diversity} is
evaluated on raw learned attention, not on the SLEEC distribution and not
on the smoothed distributions used to pool values.

For the F3 protein tower, at most 1022 residue vectors are retained. For a
longer cached sequence, the deterministic \texttt{ends\_center} policy keeps
$\lfloor1022/4\rfloor$ vectors at each end and the remaining budget from
the center, preserving their original order. The dictionary's enzyme means
are exported separately by averaging the stored residue bank; they should
not be silently replaced with the neural tower's potentially truncated mean.
EnzymeMap's specified screening library already restricts sequence length
to at most 650 amino acids.

The pooling operations scale linearly in sequence length for a fixed number
of slots. In particular, learned residue attention forms an $L\times4$
matrix, not an $L\times L$ self-attention matrix. This complexity statement
concerns the trainable pooling module after frozen features are available;
it does not describe the cost of computing ProtT5 features.

\section{Fixed chemistry descriptor}
\label{app:v4-chemistry}
The 617-dimensional chemistry vector consists of four blocks:
\begin{equation}
  617=512+64+4+37.
\end{equation}
The first block averages 512-bit radius-2 Morgan fingerprints of valid
participant molecules. The second block applies sum, mean, maximum, and
standard-deviation aggregation to 16 molecule descriptors, yielding 64
coordinates. The descriptors are exact molecular weight, heavy-atom count,
heteroatom count, formal charge, hydrogen-bond donor count, hydrogen-bond
acceptor count, topological polar surface area, logP, rotatable-bond count,
ring count, aromatic-heavy-atom fraction, stereocenter count, wildcard-atom
count, nitrogen count, oxygen count, and combined phosphorus/sulfur count.

The four set statistics are participant count, unique-participant fraction,
valid-participant fraction, and wildcard-containing fraction among valid
participants. The final block is a 37-position capacity for the recognized
core-cofactor vocabulary fitted on training reactions; unused positions are
zero. Descriptor means and standard deviations are fitted on valid training
reactions. A standard deviation below $10^{-6}$ is replaced by one.

Participant canonicalization removes atom-map indices while retaining
isomeric SMILES information. When both sides of a stored self-reaction have
the same canonical participant multiset, the participant list is used once.
Otherwise, the two sides are combined for this descriptor. Thus this chemistry
block is a participant-set feature, even when the neural molecular branches
retain physical reaction sides. Existing cofactor-related input coordinates
are distinct from the optional cofactor \emph{loss} in
Section~\ref{sec:v4-biology}; the baseline does not exclude all biological
information.

\section{Efficient evaluation of the biological objective}
\label{app:v4-efficient-bio}
For one category, define the weighted sums
\begin{equation}
  A_c=\sum_{i\in I_c}w_{ic},\qquad
  Q_c=\sum_{i\in I_c}w_{ic}^{2},\qquad
  s_c=\sum_{i\in I_c}w_{ic}\hat z_i,
\end{equation}
and the diagonal correction
$V_c=\sum_{i\in I_c}w_{ic}^{2}\|\hat z_i\|_2^2$.
The ordered off-diagonal confidence mass is $T_c=A_c^2-Q_c$. Expanding the
dot product of the weighted sum gives
\begin{equation}
  \sum_{\substack{i,j\in I_c\\i\ne j}}
      w_{ic}w_{jc}(1-\hat z_i^\top\hat z_j)
       =T_c-\|s_c\|_2^2+V_c.
  \label{eq:v4-category-sum}
\end{equation}
Retaining $V_c$ makes the expression exact even when numerical normalization
does not return exactly unit norm. For the background, define
\begin{equation}
  A_c^{\mathrm{out}}=\sum_{j\in B_c}\bar w_j^f,\qquad
  s_c^{\mathrm{out}}=\sum_{j\in B_c}\bar w_j^f\hat z_j.
\end{equation}
Then $U_c=A_c A_c^{\mathrm{out}}$ and
\begin{equation}
  \sum_{i\in I_c}\sum_{j\in B_c}
      w_{ic}\bar w_j^f(1-\hat z_i^\top\hat z_j)
       =U_c-s_c^\top s_c^{\mathrm{out}}.
  \label{eq:v4-background-sum}
\end{equation}
Background sums are obtained from the total family-annotated sum after
subtracting the category members' contributions using their \emph{family}
maximum confidence. This is important when a member's confidence for category
$c$ differs from its largest confidence in the family.

The implementation uses indexed reductions to evaluate these identities.
If $M_f$ is the number of annotated memberships in a family and $C_f$ its
number of categories, their additional time is proportional to
$(N_X+M_f+C_f)d$, for representation width $d=512$; it does not require a
dense $N_X\times N_X$ biological distance matrix. This saving does not remove
the separate $N_R\times N_E$ matrix used by the full-graph association loss.

\paragraph{Distinction from the earlier attraction-only objective.}
An earlier ablation used
\begin{equation}
  \ell_c^{\mathrm{attraction}}=
    \frac{T_c-\|s_c\|_2^2+V_c}{n_c(n_c-1)}.
\end{equation}
This encourages absolute within-category compactness. The reported
\texttt{all\_1} recipe instead uses the relative hinge in
Eq.~\eqref{eq:v4-relative-bio}. These objectives should not be interchanged
when describing or reproducing the biological experiments.

\section{Why the two association losses differ}
\label{app:v4-loss-behavior}
For a single phase-1 positive, write
$D_{ap}=\exp(z_{ap})+\sum_{n\in N_B(a)}\exp(z_{an})$. Its gradients are
\begin{align}
  \frac{\partial\ell_{\mathrm{dec}}(a,p)}{\partial z_{ap}}
      &=\frac{\exp(z_{ap})}{D_{ap}}-1,\\
  \frac{\partial\ell_{\mathrm{dec}}(a,p)}{\partial z_{an}}
      &=\frac{\exp(z_{an})}{D_{ap}},\qquad n\in N_B(a).
\end{align}
Other known positives are absent from this positive's denominator and have
zero derivative in this individual term. They receive their own positive
terms. Averaging these terms within an anchor ensures that adding more
recorded positives does not simply multiply its loss weight.

Phase 2 instead fits a categorical distribution with target probability
$1/d_a$ for each of an anchor's $d_a$ positives. Its per-logit derivative is
the predicted probability minus the target probability. Its minimum
cross-entropy for unconstrained logits approaches $\log d_a$, the entropy
of the uniform-positive target, rather than zero when $d_a>1$.
Consequently, numerical loss values from the two phases are not directly
comparable. Both losses use unobserved associations as contrastive
alternatives; neither resolves all missing-positive annotation.

\section{Annotation provenance and benchmark-specific inputs}
\label{app:v4-annotations}
The biological categories are built on the training endpoints of each target.
For EnzymeMap, EC assignments and atom mapping are read from the native
EnzymeMap cache only at indices recorded in the official training
associations. Mechanism labels are derived from the native mapped reactants
and products. The native quality field must be at least 0.5, and the retained
confidence is capped at one. Cofactor presence is recognized from the
training reaction participants and assigned confidence 0.4.

For ReactZyme, native EC annotations are restricted to training protein
sequences. Training reaction EC annotations use reaction-only Rhea matching
and exact canonical-chemistry lookup. Coarse bond-change categories use
the cached mapping of each training reaction, retain mapping confidence at
least 0.5, and multiply this confidence by the reaction-matching confidence.
Cofactor categories use the permitted training participant sets. These
reaction-only annotation resources provide information beyond the raw
released participant-set input; they are auxiliary training supervision,
not an inference-time reconstruction of missing reaction direction.

For every training edge, the reaction's observed categories are added to its
associated protein, retaining the maximum confidence when a membership is
supported more than once. No held-out reaction--enzyme association is used
for this propagation. Unknown or absent categories are omitted. In the
current EnzymeMap annotation vocabulary, the observed coarse mechanism groups
are phosphate transfer, redox/carbonyl interconversion, stereochemical
rearrangement, and sulfur/thiol chemistry. These broad labels should not be
described as a detailed model of catalytic mechanisms.

The target training graph and dictionary sizes are given in
Table~\ref{tab:v4-targets}. Duplicate association rows are removed when
building the compact phase-2 graph. Training endpoints can remain eligible
candidates under a benchmark's candidate-pool definition; the method does
not silently remove them.

\begin{table}[ht]
\centering
\small
\begin{tabular}{lrrr}
\toprule
Target & Training rows & Unique phase-2 edges & Reaction anchors\\
\midrule
ReactZyme Reaction-Sim & 147,393 & 147,393 & 6,977\\
ReactZyme Enzyme-Sim & 152,748 & 152,748 & 7,386\\
ReactZyme Time & 149,554 & 149,554 & 7,281\\
EnzymeMap & 34,427 & 34,180 & 12,603\\
\bottomrule
\end{tabular}
\caption{Target-specific training inputs for the implemented recipe. Each row
has independently fitted retrieval weights, phase-2 heads, and a dictionary.}
\label{tab:v4-targets}
\end{table}

\section{Optimization recipe and executable order}
\label{app:v4-recipe}
\begin{longtable}{p{0.50\textwidth}p{0.42\textwidth}}
\toprule
Setting & Value\\
\midrule
\endhead
Phase-1 seed; association batch size & 42; 512\\
Retained phase-1 snapshot & End of epoch 10\\
Phase-1 optimizer & AdamW, $\mathrm{lr}=10^{-4}$, weight decay $10^{-2}$\\
Phase-1 inverse temperature & Fixed $\beta=5$\\
Phase-1 directional weights & $1/2$ for each direction\\
Unknown-negative weight & 1\\
Entropy and diversity coefficients & 0.01 and 0.01\\
Protein attention uniform mixture & $\eta=0.05$\\
Gate uniform mixture & $\zeta=0.05$\\
Phase-2 optimizer & AdamW, $\mathrm{lr}=10^{-4}$, weight decay $10^{-3}$\\
Phase-2 updates; seed & 100 full-graph updates; 42\\
Phase-2 temperature; identity coefficient & $\tau_2=0.2$; 10\\
Phase-2 residual training coefficient & $\kappa=0.2$\\
Phase-2 gradient clipping & Global gradient norm 1\\
Biological mode; margin & Relative; $m=0.1$\\
Biological family coefficients, \texttt{all\_1} & $(\lambda_{\mathrm{EC}},\lambda_{\mathrm{cof}},\lambda_{\mathrm{mech}})=(1,1,1)$\\
Biological family divisor & Fixed at 3, including removal controls\\
Phase-1 biological supervision, \texttt{all\_1} & None\\
Inference protein fusion multiplier & $\nu=3$\\
Inference residual multiplier & $c=0.5$\\
Dictionary score coefficient & $\alpha=0.4$\\
Protein dictionary neighbors & 32\\
Dictionary kernel temperatures & $\tau_E=\tau_R=0.03$\\
Reaction dictionary truncation & None; all training reaction anchors\\
\bottomrule
\end{longtable}

\paragraph{Training and inference procedure.}
\begin{enumerate}[leftmargin=*,itemsep=5pt]
  \item Construct the target training association graph and the frozen feature
  caches. Initialize the trainable retrieval components for that target and
  load the frozen pretrained feature resources and SLEEC scorer.
  \item Train the phase-1 encoders with Eq.~\eqref{eq:v4-phase1-total}. For
  each batch, deduplicate endpoint IDs and recover all known in-batch training
  positives. Retain the epoch-10 state.
  \item Freeze F3. Export native training embeddings using its learned scale
  $s$, without the factor-three inference adjustment. Build the dictionary
  centers, training protein bank, training reaction bank, and association map.
  \item Initialize the residual heads at identity. For each of 100 updates,
  encode every training endpoint, evaluate the full bidirectional
  cross-entropy and identity loss, and optionally add the biological relative
  loss. Clip gradients and update only the residual heads.
  \item At inference, obtain the protein's global and fused branches and use
  $\nu s$ in Eq.~\eqref{eq:v4-inference-fusion}. Apply the trained heads with
  residual coefficient $c\kappa=0.1$.
  \item Independently compute the enzyme and reaction dictionary coordinates
  from frozen raw features and the stored training dictionary. Concatenate
  the weighted blocks, or retain equivalent dense and sparse blocks.
  \item Rank candidates by Eq.~\eqref{eq:v4-final-score} in the requested
  direction. No candidate labels, candidate-pool normalization, or biological
  annotation files are needed for encoding new endpoints.
\end{enumerate}

\paragraph{Checkpoint representation of the fusion adjustment.}
The EnzymeMap deployment checkpoint already stores the factor-three change
to the bounded fusion scalar, so its runtime fusion multiplier is one.
The ReactZyme deployment checkpoints store the native scalar and apply three
at runtime. These represent the same recipe. Applying three again to the
calibrated EnzymeMap checkpoint would implement a different model. Phase-2
training used the native feature exports for every target.

\paragraph{Numerical evaluation.}
Dictionary feature norms, centers, and nearest-neighbor similarities use
FP64 reductions where implemented before storing FP32 features. Neighbor
selection resolves cutoff ties by the lower training-anchor index.
Equivalent mathematical score decompositions can nevertheless have different
floating-point rounding when executed in different orders. The frozen
evaluation routines and candidate ordering should therefore be retained,
particularly for nearly tied ReactZyme reaction candidates. The dictionary
is part of the deployed method, not a post-hoc normalization fitted to the
current screening library.

\section{Scope of the biological extension}
\label{app:v4-scope}
The main specification above is the phase-2-only biological configuration.
Experiments that retrain F3 with biological supervision are a different
variant: they add the family-averaged loss to the phase-1 objective, using
deduplicated minibatch endpoints rather than the full graph. In that case
the gradient can reach the learned residue queries and reaction encoder.
Category eligibility and background coverage then depend on each minibatch,
so this is not equivalent to adding the same nominal coefficient only in
phase 2.

The architecture and loss definition describe how biological relationships
are introduced, not whether they improve generalization. That empirical
question requires matched controls, including removal and shuffled-label
experiments. Equal family coefficients do not imply equal gradient influence
or equal predictive benefit. The 512 neural coordinates have no fixed
biochemical names; only the dictionary coordinates have an explicit
training-reaction interpretation.
